% ---------------------------------------------------------------
\section{Conclusion and Future Work}
\label{sec:conclusion}
This paper presented \idm{}, a label-safe monitoring framework that keeps label-free
input evidence separate from calibrated impact evidence and scores both against harm
that is verified only after a release delay. Across controlled, replicated and natural
streams, the framework detected marginal change reliably and estimated accuracy well
while the source correctness relation held. It also showed where unlabelled monitoring
cannot work. Conditional-only change caused 1,205 harmful batches that were never
flagged; nuisance change raised alarms without harm; and calibration failed to
transfer on drifting sensors. A dependence statistic closed one detection blind spot, but
combined with KS it nearly tripled the stationary alarm burden. Released labels repaired estimation on one natural stream but
not on another. Monitoring therefore cannot be judged by its alarms alone. It should be
judged by verified harm, with the release delay treated as part of the design.

Future work should evaluate the framework on more independent natural streams with
recorded label-arrival times. It should design label-acquisition policies with known
inclusion probabilities and quantify their costs, and propagate calibration and
reference uncertainty into the interval. Detector thresholds should be calibrated with
the stationary alarm burden as an explicit target, and a pre-registered analyst study
should test whether the aligned view changes decisions.
